% !TEX root = ../algebra_02.tex

\section{Собственные значения и собственные векторы}

\begin{Definition}{Собственное значение и собственный вектор}{}
	Пусть $\alpha\in\End V$.
	Число $\lambda\in K$ называется собственным значением оператора $\alpha$, если существует ненулевой вектор $v\in V$ такой, что
	\[
	\alpha v=\lambda v.
	\]
	Такой вектор $v$ называется собственным вектором, отвечающим собственному значению $\lambda$.
\end{Definition}

\begin{Definition}{Собственное подпространство}{}
	Собственным подпространством, отвечающим $\lambda$, называется
	\[
	V_\lambda=\{v\in V\mid \alpha v=\lambda v\}.
	\]
	Эквивалентно,
	\[
	V_\lambda=\Ker(\alpha-\lambda\operatorname{id}_V).
	\]
\end{Definition}

\begin{Definition}{Геометрическая кратность}{}
	Если $\lambda$ --- собственное значение оператора $\alpha$, то число
	\[
	g_\lambda:=\dim V_\lambda
	\]
	называется геометрической кратностью собственного значения $\lambda$.
\end{Definition}

\begin{Theorem}{Собственные векторы с разными собственными значениями}{}
	Пусть $\lambda_1,\ldots,\lambda_k$ --- попарно различные собственные значения оператора $\alpha$.
	Если $v_i\in V_{\lambda_i}$ и $v_i\ne 0$, то векторы
	$v_1,\ldots,v_k$ линейно независимы.
\end{Theorem}

\begin{proof}
	Докажем индукцией по $k$.
	
	При $k=1$ утверждение очевидно, так как $v_1\ne 0$.
	
	Пусть
	\[
	a_1v_1+\cdots+a_kv_k=0.
	\]
	Применим $\alpha$:
	\[
	a_1\lambda_1v_1+\cdots+a_k\lambda_kv_k=0.
	\]
	Умножим исходное равенство на $\lambda_k$ и вычтем его из последнего:
	\[
	a_1(\lambda_1-\lambda_k)v_1+\cdots+
	a_{k-1}(\lambda_{k-1}-\lambda_k)v_{k-1}=0.
	\]
	По предположению индукции $v_1,\ldots,v_{k-1}$ линейно независимы.
	Так как $\lambda_i\ne\lambda_k$ при $i<k$, получаем
	$a_1=\cdots=a_{k-1}=0$.
	Тогда из исходного равенства следует $a_kv_k=0$, откуда $a_k=0$.
\end{proof}

\begin{Corollary}{Прямая сумма собственных подпространств}{}
	Если $\lambda_1,\ldots,\lambda_k$ --- различные собственные значения $\alpha$, то
	\[
	V_{\lambda_1}+\cdots+V_{\lambda_k}
	=
	V_{\lambda_1}\oplus\cdots\oplus V_{\lambda_k}.
	\]
\end{Corollary}

\begin{Corollary}{Оценка числа собственных значений}{}
	Пусть $\lambda_1,\ldots,\lambda_k$ --- все различные собственные значения оператора $\alpha$.
	Тогда
	\[
	g_{\lambda_1}+\cdots+g_{\lambda_k}\le \dim V.
	\]
	В частности, у оператора на $n$-мерном пространстве не больше $n$ различных собственных значений.
\end{Corollary}

\begin{Corollary}{Достаточное условие диагонализируемости}{}
	Если $\dim V=n$ и у оператора $\alpha$ есть $n$ различных собственных значений, то $\alpha$ диагонализируем.
\end{Corollary}

\begin{proof}
	Выберем ненулевой собственный вектор $v_i\in V_{\lambda_i}$ для каждого собственного значения.
	По теореме эти $n$ векторов линейно независимы.
	Значит, они образуют базис $V$, и в этом базисе матрица оператора диагональна.
\end{proof}

\subsection{Диагонализируемые операторы}

\begin{Definition}{Диагонализируемый оператор}{}
	Оператор $\alpha\in\End V$ называется диагонализируемым, если существует базис $E$ пространства $V$, в котором матрица $[\alpha]_E$ диагональна.
\end{Definition}

\begin{Theorem}{Критерий диагонализируемости}{}
	Пусть $\lambda_1,\ldots,\lambda_k$ --- все различные собственные значения оператора $\alpha$.
	Тогда $\alpha$ диагонализируем тогда и только тогда, когда
	\[
	V=V_{\lambda_1}\oplus\cdots\oplus V_{\lambda_k}.
	\]
\end{Theorem}

\begin{proof}
	Если $\alpha$ диагонализируем, то существует базис из собственных векторов.
	Группируя эти векторы по собственным значениям, получаем разложение $V$ в прямую сумму собственных подпространств.
	
	Обратно, если
	$V=V_{\lambda_1}\oplus\cdots\oplus V_{\lambda_k}$,
	то в каждом $V_{\lambda_i}$ выберем базис.
	Объединение этих базисов даёт базис $V$, состоящий из собственных векторов.
	В таком базисе матрица оператора диагональна.
\end{proof}

\begin{Lemma}{Собственные значения диагонального оператора}{}
	Пусть $E=(e_1,\ldots,e_n)$ --- базис $V$ и
	\[
	[\alpha]_E=
	\operatorname{diag}(
	\underbrace{\lambda_1,\ldots,\lambda_1}_{m_1},
	\underbrace{\lambda_2,\ldots,\lambda_2}_{m_2},
	\ldots,
	\underbrace{\lambda_k,\ldots,\lambda_k}_{m_k}
	),
	\]
	где $\lambda_1,\ldots,\lambda_k$ попарно различны.
	
	Тогда $\lambda_1,\ldots,\lambda_k$ --- все собственные значения $\alpha$, причём
	\[
	g_{\lambda_i}=m_i
	\qquad
	(i=1,\ldots,k).
	\]
\end{Lemma}

\begin{proof}
	Для любого $\lambda\in K$
	\[
	[\alpha-\lambda\operatorname{id}_V]_E
	=
	[\alpha]_E-\lambda I_n.
	\]
	Если $\lambda\notin\{\lambda_1,\ldots,\lambda_k\}$, то все диагональные элементы этой матрицы ненулевые, значит,
	$\Ker(\alpha-\lambda\operatorname{id}_V)=0$.
	Следовательно, такое $\lambda$ не является собственным значением.
	
	Пусть теперь $\lambda=\lambda_i$.
	Тогда в матрице $[\alpha-\lambda_i\operatorname{id}_V]_E$ нули стоят ровно на тех местах диагонали, которые соответствуют блоку $\lambda_i$.
	Поэтому
	\[
	V_{\lambda_i}
	=
	\Lin(e_{m_1+\cdots+m_{i-1}+1},\ldots,e_{m_1+\cdots+m_i}),
	\]
	и значит, $\dim V_{\lambda_i}=m_i$.
\end{proof}


\begin{Theorem}{Критерий через геометрические кратности}{}
	Оператор $\alpha$ диагонализируем тогда и только тогда, когда
	\[
	g_{\lambda_1}+\cdots+g_{\lambda_k}=n.
	\]
\end{Theorem}

\begin{proof}
	Пусть $\alpha$ диагонализируем. Тогда существует базис $E$ из собственных векторов.
	Разобьём этот базис на группы по собственным значениям. Если собственному значению
	$\lambda_i$ соответствует $m_i$ векторов базиса, то по лемме о диагональном операторе
	$g_{\lambda_i}=m_i$. Поэтому
	\[
	g_{\lambda_1}+\cdots+g_{\lambda_k}
	=
	m_1+\cdots+m_k
	=
	n.
	\]
	
	Обратно, пусть
	$g_{\lambda_1}+\cdots+g_{\lambda_k}=n$.
	Собственные подпространства, отвечающие разным собственным значениям, образуют прямую сумму:
	\[
	V_{\lambda_1}+\cdots+V_{\lambda_k}
	=
	V_{\lambda_1}\oplus\cdots\oplus V_{\lambda_k}.
	\]
	Размерность этой суммы равна $n$, значит,
	\[
	V=V_{\lambda_1}\oplus\cdots\oplus V_{\lambda_k}.
	\]
	Выберем базис в каждом $V_{\lambda_i}$. Объединение этих базисов даёт базис $V$,
	состоящий из собственных векторов. Значит, $\alpha$ диагонализируем.
\end{proof}

\subsection{Характеристический многочлен}

Пусть $\alpha\in \operatorname{End} V$, $\dim V=n$, и $E$ --- базис $V$.
Обозначим
\[
A=[\alpha]_E.
\]

\begin{Statement}{Собственные значения как корни определителя}{}
	Число $\lambda\in K$ является собственным значением оператора $\alpha$ тогда и только тогда, когда
	\[
	\det(A-\lambda I_n)=0.
	\]
\end{Statement}

\begin{proof}
	Имеем цепочку эквивалентностей:
	\[
	\lambda \text{ --- собственное значение}
	\Longleftrightarrow
	\ker(\alpha-\lambda\operatorname{id}_V)\ne 0
	\Longleftrightarrow
	\det[\alpha-\lambda\operatorname{id}_V]_E=0.
	\]
	Но
	\[
	[\alpha-\lambda\operatorname{id}_V]_E=A-\lambda I_n.
	\]
	Следовательно, $\lambda$ является собственным значением тогда и только тогда, когда
	$\det(A-\lambda I_n)=0$.
\end{proof}

\begin{Definition}{Характеристический многочлен матрицы}{}
	Пусть $A\in M_n(K)$. Характеристическим многочленом матрицы $A$ называется
	\[
	\chi_A(x)=\det(A-xI_n).
	\]
\end{Definition}

\begin{Remark}{Монический вариант}{}
	Так как $\chi_A(x)=\det(A-xI_n)$ имеет старший коэффициент $(-1)^n$, часто удобно рассматривать
	монический многочлен
	\[
	p_A(x)=(-1)^n\chi_A(x)=\det(xI_n-A).
	\]
	Корни и их кратности у $\chi_A$ и $p_A$ одинаковые.
\end{Remark}

\begin{Theorem}{Старшие коэффициенты характеристического многочлена}{}
	Пусть $A=(a_{ij})\in M_n(K)$. Тогда
	\[
	(-1)^n\chi_A(x)
	=
	x^n-\operatorname{tr}(A)x^{n-1}+\cdots+(-1)^n\det A,
	\]
	где
	\[
	\operatorname{tr}(A)=a_{11}+\cdots+a_{nn}.
	\]
\end{Theorem}

\begin{proof}
	Рассмотрим
	\[
	(-1)^n\chi_A(x)=\det(xI_n-A).
	\]
	В формуле для определителя старшая степень $x^n$ получается только из произведения
	диагональных элементов:
	\[
	(x-a_{11})(x-a_{22})\cdots(x-a_{nn}).
	\]
	Поэтому старший коэффициент равен $1$.
	
	Коэффициент при $x^{n-1}$ получается из той же диагональной перестановки, когда один раз
	берётся $-a_{ii}$, а остальные разы берётся $x$. Поэтому этот коэффициент равен
	\[
	-(a_{11}+\cdots+a_{nn})=-\operatorname{tr}(A).
	\]
	
	Свободный член равен значению при $x=0$:
	\[
	\det(0\cdot I_n-A)=\det(-A)=(-1)^n\det A.
	\]
	Отсюда получается указанная формула.
\end{proof}

\begin{Theorem}{Независимость характеристического многочлена от базиса}{}
	Пусть $E$ и $E'$ --- базисы пространства $V$, $\alpha\in\operatorname{End} V$.
	Тогда
	\[
	\chi_{[\alpha]_E}=\chi_{[\alpha]_{E'}}.
	\]
\end{Theorem}

\begin{proof}
	Пусть $A=[\alpha]_E$, $A'=[\alpha]_{E'}$.
	Тогда матрицы $A$ и $A'$ подобны, то есть существует обратимая матрица $C$ такая, что
	\[
	A'=C^{-1}AC.
	\]
	Следовательно,
	\begin{align*}
		\chi_{A'}(x)
		&=\det(A'-xI_n)\\
		&=\det(C^{-1}AC-xI_n)\\
		&=\det(C^{-1}(A-xI_n)C)\\
		&=\det(C^{-1})\det(A-xI_n)\det(C)\\
		&=\det(A-xI_n)
		=
		\chi_A(x).
	\end{align*}
\end{proof}

\begin{Definition}{Характеристический многочлен оператора}{}
	Характеристическим многочленом оператора $\alpha\in\operatorname{End} V$ называется
	\[
	\chi_\alpha(x):=\chi_{[\alpha]_E}(x),
	\]
	где $E$ --- любой базис $V$.
\end{Definition}

\begin{Corollary}{Собственные значения оператора}{}
	Число $\lambda\in K$ является собственным значением оператора $\alpha$ тогда и только тогда, когда
	\[
	\chi_\alpha(\lambda)=0.
	\]
\end{Corollary}

\subsection{Характеристический многочлен и инвариантные подпространства}

\begin{Theorem}{Факторизация по инвариантному подпространству}{}
	Пусть $\alpha\in\operatorname{End} V$, а $W\le V$ --- инвариантное подпространство.
	Тогда
	\[
	\chi_\alpha(x)
	=
	\chi_{\alpha|_W}(x)\chi_{\overline{\alpha}}(x),
	\]
	где $\overline{\alpha}$ --- индуцированный оператор на $V/W$.
\end{Theorem}

\begin{proof}
	Пусть $\dim W=m$.
	Возьмём базис $e_1,\ldots,e_m$ пространства $W$ и дополним его до базиса
	\[
	E=(e_1,\ldots,e_m,e_{m+1},\ldots,e_n)
	\]
	пространства $V$.
	
	Так как $W$ инвариантно, матрица оператора $\alpha$ в базисе $E$ имеет блочный вид
	\[
	[\alpha]_E=
	\begin{pmatrix}
		B & *\\
		0 & C
	\end{pmatrix},
	\]
	где $B=[\alpha|_W]$ в базисе $e_1,\ldots,e_m$, а $C=[\overline{\alpha}]$ в базисе
	$\overline e_{m+1},\ldots,\overline e_n$ пространства $V/W$.
	
	Тогда
	\[
	[\alpha]_E-xI_n=
	\begin{pmatrix}
		B-xI_m & *\\
		0 & C-xI_{n-m}
	\end{pmatrix}.
	\]
	Определитель блочно-треугольной матрицы равен произведению определителей диагональных блоков:
	\[
	\chi_\alpha(x)
	=
	\det(B-xI_m)\det(C-xI_{n-m})
	=
	\chi_{\alpha|_W}(x)\chi_{\overline{\alpha}}(x).
	\]
\end{proof}

\begin{Corollary}{Случай прямой суммы инвариантных подпространств}{}
	Если
	\[
	V=W_1\oplus W_2,
	\]
	причём $W_1$ и $W_2$ инвариантны относительно $\alpha$, то
	\[
	\chi_\alpha(x)
	=
	\chi_{\alpha|_{W_1}}(x)\chi_{\alpha|_{W_2}}(x).
	\]
\end{Corollary}

\subsection{Алгебраическая кратность}

\begin{Definition}{Алгебраическая кратность}{}
	Пусть $\lambda$ --- собственное значение оператора $\alpha$.
	Алгебраической кратностью $\lambda$ называется кратность $\lambda$ как корня
	характеристического многочлена $\chi_\alpha$.
	
	Обозначение:
	\[
	a_\lambda.
	\]
\end{Definition}

\begin{Theorem}{Связь геометрической и алгебраической кратностей}{}
	Для любого собственного значения $\lambda$ оператора $\alpha$ выполнено
	\[
	1\le g_\lambda\le a_\lambda.
	\]
\end{Theorem}

\begin{proof}
	Пусть $W=V_\lambda$.
	Тогда $W$ инвариантно относительно $\alpha$, поскольку для любого $v\in V_\lambda$
	имеем $\alpha v=\lambda v\in V_\lambda$.
	
	Применим теорему о факторизации характеристического многочлена:
	\[
	\chi_\alpha(x)
	=
	\chi_{\alpha|_{V_\lambda}}(x)\chi_{\overline{\alpha}}(x).
	\]
	На $V_\lambda$ оператор $\alpha$ действует как умножение на $\lambda$, поэтому
	\[
	[\alpha|_{V_\lambda}]=\lambda I_{g_\lambda}.
	\]
	Следовательно,
	\[
	\chi_{\alpha|_{V_\lambda}}(x)
	=
	\det(\lambda I_{g_\lambda}-xI_{g_\lambda})
	=
	(\lambda-x)^{g_\lambda}.
	\]
	Значит, $(\lambda-x)^{g_\lambda}$ делит $\chi_\alpha(x)$.
	Отсюда $g_\lambda\le a_\lambda$.
	Неравенство $g_\lambda\ge 1$ следует из того, что $\lambda$ --- собственное значение.
\end{proof}

\subsection{Критерий диагонализируемости через характеристический многочлен}

\begin{Theorem}{Критерий через алгебраические кратности}{}
	Пусть $\alpha\in\operatorname{End} V$, $\dim V=n$.
	Оператор $\alpha$ диагонализируем тогда и только тогда, когда выполнены два условия:
	
	\begin{enumerate}
		\item характеристический многочлен $\chi_\alpha$ раскладывается над $K$ на линейные множители;
		\item для каждого собственного значения $\lambda$ выполнено $g_\lambda=a_\lambda$.
	\end{enumerate}
\end{Theorem}

\begin{proof}
	Пусть $\alpha$ диагонализируем. Тогда в некотором базисе его матрица диагональна:
	\[
	[\alpha]_E=
	\operatorname{diag}(
	\underbrace{\lambda_1,\ldots,\lambda_1}_{m_1},
	\underbrace{\lambda_2,\ldots,\lambda_2}_{m_2},
	\ldots,
	\underbrace{\lambda_k,\ldots,\lambda_k}_{m_k}
	),
	\]
	где $\lambda_1,\ldots,\lambda_k$ --- различные собственные значения.
	
	Тогда
	\[
	\chi_\alpha(x)
	=
	(\lambda_1-x)^{m_1}\cdots(\lambda_k-x)^{m_k}.
	\]
	Значит, характеристический многочлен раскладывается на линейные множители.
	Кроме того, по лемме о диагональном операторе $g_{\lambda_i}=m_i$.
	С другой стороны, $a_{\lambda_i}=m_i$ по виду характеристического многочлена.
	Следовательно, $g_{\lambda_i}=a_{\lambda_i}$ для всех $i$.
	
	Обратно, пусть $\chi_\alpha$ раскладывается на линейные множители и
	$g_\lambda=a_\lambda$ для каждого собственного значения $\lambda$.
	Так как корни характеристического многочлена --- это ровно собственные значения, получаем
	\[
	n=\sum_{\lambda} a_\lambda.
	\]
	По условию $a_\lambda=g_\lambda$, значит,
	\[
	n=\sum_{\lambda} g_\lambda.
	\]
	По критерию через геометрические кратности оператор $\alpha$ диагонализируем.
\end{proof}

\begin{Corollary}{Случай простых корней}{}
	Если характеристический многочлен $\chi_\alpha$ имеет $n$ различных корней в поле $K$,
	то оператор $\alpha$ диагонализируем.
\end{Corollary}

\begin{proof}
	В этом случае все алгебраические кратности равны $1$.
	Для каждого собственного значения $\lambda$ имеем
	\[
	1\le g_\lambda\le a_\lambda=1,
	\]
	поэтому $g_\lambda=a_\lambda=1$.
	Значит, выполнен критерий диагонализируемости через алгебраические кратности.
\end{proof}